\section{The ACDB Benchmark}
\label{sec:benchmark}

\textbf{ACDB} is a push-based benchmark environment in which an agent receives observational data, spends a budget of hard single-node interventions, and submits an estimated causal graph. A benchmark instance is $(G, \mathcal{M}, B, D^{\mathrm{obs}})$: a ground-truth DAG, a linear--Gaussian SCM, an intervention budget, and a fixed observational sample. The environment hides $G$ and $\mathcal{M}$ and exposes only $D^{\mathrm{obs}}$, $d$, and $B$.

\subsection{World generation}

Each instance assumes causal sufficiency, linear--Gaussian noise (Eq.~\ref{eq:linear-gaussian}), and edge weights drawn uniformly from $[-2,-0.5]\cup[0.5,2]$. We sample a uniform topological order, select exactly $k$ forward edges, and randomly permute node indices so names do not leak topology. Instances are rejected if any relevant partial correlation falls below $\texttt{faithfulness\_eps}=0.1$. For each accepted instance we precompute the minimum single-node intervention set $\mathcal{I}^\star$ that orients every CPDAG-undirected edge; this defines the efficiency baseline.

\subsection{Session API}

An ACDB session enforces a strict observe--intervene--submit protocol. \texttt{observe()} returns the pre-sampled $D^{\mathrm{obs}}$ once. \texttt{intervene(var, value)} draws $n_{\mathrm{int}}$ fresh samples under the hard intervention $\mathrm{do}(X_{\texttt{var}} = \texttt{value})$ and decrements the budget; requests beyond $B$ raise. \texttt{submit\_graph(submission)} seals the session. Submissions are mixed graphs of directed and undirected edges; undirected edges represent orientations the agent was unable or unwilling to commit to. Interventions are hard (the target variable is set exogenously; its parents are severed during sampling), and the scoring contract treats an unresolved edge as a distinct failure mode from a reversal or omission.

\subsection{Scoring contract}
\label{sec:scoring}

ACDB reports three complementary views of a submission; all are deterministic given $(G,\text{submission},B)$.

\paragraph{Observational layer: skeleton and compelled F1.} Against the true CPDAG of $G$, we compute standard precision/recall/F1 on the undirected skeleton ($\mathit{skeleton\_f1}$) and on the compelled (observationally orientable) edges ($\mathit{compelled\_f1}$). A reversed compelled edge counts as both a false positive on the reversed direction and a false negative on the true direction, so the agent is not rewarded for ``at least the adjacency was right'' when it flipped a compelled orientation.

\paragraph{DAG layer: directed F1 and SHD.} Against the true DAG $G$ itself, we compute set-intersection precision/recall/F1 on directed edges ($\mathit{directed\_f1}$) and a structural Hamming distance $\mathit{dag\_shd}$. SHD counts, over the union of true and submitted adjacencies, the following errors at one per edge: \emph{extra} (submission has an edge, truth does not), \emph{missing} (truth has an edge, submission has nothing), \emph{reversed} (both have a directed edge but in opposite directions), and \emph{unresolved} (truth has a directed edge, submission has an undirected edge on the same pair). Note that a reversal is one error, not two.

\paragraph{Efficiency layer.} Given the precomputed optimal intervention set $\mathcal{I}^\star$ of size $|\mathcal{I}^\star|$ and the number of interventions actually used, the efficiency score is
\begin{equation}
\eta \;=\; \frac{|\mathcal{I}^\star|}{\max(|\mathcal{I}_{\mathrm{used}}|, |\mathcal{I}^\star|)}.
\label{eq:efficiency}
\end{equation}
This is one whenever the agent used at most the optimal number of interventions (including zero-budget submission when $\mathcal{I}^\star = \varnothing$), and decreases toward zero as wasted interventions accumulate.

\paragraph{What the scoring contract deliberately excludes.} There is no MEC-aware partial credit in the DAG layer: the observational layer already expresses MEC-ceiling performance, and making the DAG layer MEC-aware would collapse the two.

\subsection{The difficulty ladder}

ACDB organizes instances into six levels, summarized in Table~\ref{tab:ladder}. Each level is a triple of difficulty axes: graph size and density, statistical power (sample counts and noise variance), and budget tightness (\texttt{budget\_slack} controls how many interventions above $|\mathcal{I}^\star|$ are allotted). The ladder is not a single monotone axis; levels 2 and 4 are designed to hold graph size fixed and isolate one axis at a time (statistical in L2, budget in L4), while L3 and L5 are structural stress levels.

\begin{table}[t]
\centering
\caption{ACDB ladder. Each level fixes $d$ (variables), $k$ (true edge count), $n_{\mathrm{obs}}$ / $n_{\mathrm{int}}$ (sample counts), $\sigma^2$ (noise variance), and \texttt{slack} (interventions above the precomputed optimal). The intervention budget is $B = |\mathcal{I}^\star| + \texttt{slack}$ and is therefore instance-specific.}
\label{tab:ladder}
\begin{tabular}{cllcccccc}
\toprule
Level & Role & $d$ & $k$ & $n_{\mathrm{obs}}$ & $n_{\mathrm{int}}$ & $\sigma^2$ & slack \\
\midrule
0 & tutorial       & 4 & 5 & 50 & 25 & 0.5 & 2 \\
1 & standard       & 5 & 6 & 25 & 15 & 1.0 & 1 \\
2 & statistical    & 5 & 6 & 15 & 10 & 1.5 & 1 \\
3 & structural     & 7 & 9 & 25 & 15 & 1.0 & 1 \\
4 & pressure       & 5 & 6 & 25 & 15 & 1.0 & 0 \\
5 & hard           & 7 & 9 & 15 & 10 & 1.5 & 0 \\
\bottomrule
\end{tabular}
\end{table}

Each level is evaluated on 8 seeds, preflight-validated so that instance assembly (DAG sample $\to$ faithfulness check $\to$ SCM parameterization $\to$ optimal intervention set) succeeds for all seeds in the manifest. The same manifest is used across all agents and baselines for a given ladder run, so comparisons are paired.

\subsection{Density leakage: a benchmark-design hazard}
\label{sec:density}

A structure-blind random baseline --- one that samples a random DAG using only the public variable count $d$, with no access to the true edge count $k$ --- provides a floor for what any non-trivial agent must beat. When graph density $\rho = k / \binom{d}{2}$ is high, this floor is also high: a randomly sampled DAG agrees with the truth on a non-trivial fraction of directed edges simply because both contain many edges. We call this effect \emph{density leakage}.

Empirically, on the ladder of Table~\ref{tab:ladder}, the structure-blind random baseline attains overall directed F1 of $0.236$ across $4800$ samples (Section~\ref{sec:results}). Measuring the same baseline at several post hoc densities, without giving the baseline access to $k$, we observe an approximately monotone heuristic:
\begin{equation}
\mathbb{E}[\text{Dir-F1}_{\text{random}}] \;\approx\; \frac{\rho}{1 + 2\rho}, \qquad \rho = \frac{k}{\binom{d}{2}}.
\label{eq:density-leakage}
\end{equation}
Equation~\eqref{eq:density-leakage} is a probe fit, not a derived bound. The density-probe ladder ($\rho \le 0.33$ on non-tutorial levels) reduces the measured random floor from $23.6\%$ to $16.9\%$ overall. We keep the original ladder for v0 because the LLM and PC runs were collected there, and use the probe to define the v1 calibration target.

% [FIGURE: three-tier scoring contract diagram] -- placeholder for v0.
% [FIGURE: random baseline Dir F1 vs density rho per ladder level] -- placeholder for v0.
